\rhead{PX101: Physics Lab}
\phantomsection 
\section*{Physics Lab (PX101)}
\label{course:PX101}

\noindent \small{\hyperref[main:scheme]{$\leftarrow$ Back to Scheme}} \vspace{0.5cm}

\noindent \textbf{Credits:} 2 (0-0-2)

\subsection*{Course Content}
\noindent\textbf{Lab 1: Semiconductor Diodes} \
Focuses on analyzing the V-I characteristics of a P-N Junction Diode in forward and reverse bias to understand the fundamental building block of logic gates and rectification. \\

\vspace{0.2cm}

\noindent\textbf{Lab 2: Voltage Regulation} \
Focuses on the breakdown mechanism of a Zener Diode and its application as a Voltage Regulator, simulating power stability in computer power supply units (PSUs). \\

\vspace{0.2cm}

\noindent\textbf{Lab 3: Optoelectronics and Sensors} \
Focuses on the characteristics of a Solar Cell or Photodiode, understanding how light energy is converted to electrical signals, applicable to optical sensors and energy harvesting. \\

\vspace{0.2cm}

\noindent\textbf{Lab 4: Energy Band Gap} \
Focuses on determining the Energy Band Gap of a semiconductor material using the Four-Probe Method, linking temperature changes to conductivity (hardware thermal management). \\

\vspace{0.2cm}

\noindent\textbf{Lab 5: Quantum Constants} \
Focuses on the determination of Planck’s Constant using LEDs of different wavelengths, verifying the quantum mechanical relationship between energy and frequency. \\

\vspace{0.2cm}

\noindent\textbf{Lab 6: Optical Fiber Communication} \
Focuses on determining the Numerical Aperture (NA) and bending losses of an Optical Fiber, simulating data transmission efficiency and signal attenuation in network cables. \\

\vspace{0.2cm}

\noindent\textbf{Lab 7: Wave Optics and Diffraction} \
Focuses on determining the wavelength of a Laser source using a Diffraction Grating, illustrating wave interference principles used in holography and optical storage. \\

\vspace{0.2cm}

\noindent\textbf{Lab 8: Electromagnetism} \
Focuses on the variation of the Magnetic Field along the axis of a circular coil (Stewart-Gee’s Experiment), visualizing field lines and electromagnetic induction principles. \\

\vspace{0.2cm}

\noindent\textbf{Lab 9: Resonance and Signals} \
Focuses on the study of standing waves and resonance frequency using Melde’s Experiment, providing a physical analogy for signal processing and harmonic analysis. \\

\vspace{0.2cm}

\noindent\textbf{Lab 10: Rigid Body Dynamics} \
Focuses on determining the acceleration due to gravity ($g$) and Radius of Gyration using a Compound Pendulum, verifying physical laws used in rigid body physics engines.

\vspace{0.2cm}

\newpage
\rhead{Curriculum Schema}